\documentclass{article}

\usepackage{amsmath}
\usepackage{amssymb}
\usepackage{polski}

\usepackage[utf8]{inputenc}
\usepackage[T1]{fontenc}

\usepackage{tikz}

\usepackage{listings}
\lstset{
    language=SQL,
    inputencoding=utf8x,
    extendedchars=\true,
    literate={ą}{{\k{a}}}1
             {Ą}{{\k{A}}}1
             {ę}{{\k{e}}}1
             {Ę}{{\k{E}}}1
             {ó}{{\'o}}1
             {Ó}{{\'O}}1
             {ś}{{\'s}}1
             {Ś}{{\'S}}1
             {ł}{{\l{}}}1
             {Ł}{{\L{}}}1
             {ż}{{\.z}}1
             {Ż}{{\.Z}}1
             {ź}{{\'z}}1
             {Ź}{{\'Z}}1
             {ć}{{\'c}}1
             {Ć}{{\'C}}1
             {ń}{{\'n}}1
             {Ń}{{\'N}}1
}

\newcommand\tikzmark[1]{%
\tikz[remember picture,baseline] \node[inner sep=0,outer sep=0] (#1){\rule{0em}{1.75em}};%
}

\newcommand\Overline[3][line width=1pt]{%
    \begin{tikzpicture}[remember picture,overlay]
      \draw[#1] (#2.north east) -- (#3.north west);
    \end{tikzpicture}
    }

\title{Liczby zespolone - zestaw II}
\date{2024}
% \author{Maciej Różewicz}

\begin{document}
\maketitle

\begin{enumerate}
    \item Przedstawić w postaci wykładniczej:
    \begin{enumerate}
        \item $-1$,
        \item $1+i$,
        \item $-i$,
        \item $1-i\sqrt{3}$,
        \item $-2+7i$
    \end{enumerate}

    \item Wykonać działania:
    \begin{enumerate}
        \item $2e^{i\frac{\pi}{4}} 3e^{i\frac{\pi}{2}}$,
        \item $\frac{2e^{i\frac{\pi}{4}}}{3e^{i\frac{\pi}{2}}}$,
        \item $-(2e^{i\frac{\pi}{4}})$,
        \item $(2e^{i\frac{\pi}{4}})^{5}$,
        \item $\overline{(2e^{i\frac{\pi}{4}})}$,
        \item $\frac{1}{2e^{i\frac{\pi}{4}}}$
    \end{enumerate}

    \item Wyznaczyć pierwiastki liczby zespolonej:
    \begin{enumerate}
        \item $\sqrt{2i}$,
        \item $\sqrt[3]{i}$,
        \item $\sqrt{-4}$,
        \item $\sqrt[4]{1}$,
        \item $\sqrt[3]{8i}$,
        \item $\sqrt[4]{-\frac{1}{2} + \frac{\sqrt{3}}{2}i}$,
        \item $\sqrt[8]{1}$
    \end{enumerate}
    
    \item Rozwiązać równanie w dziedzinie liczb zespolonych:
    \begin{enumerate}
        \item $z^{2} + 3z + 1 = 0$,
        \item $z^{2} + 4z + 5 = 0$,
        \item $z^{2} - \sqrt{2}(1+i) -\frac{3}{2}i = 0$,
        \item $z^{3} + 8 = 0$
    \end{enumerate}

\end{enumerate}

\newpage
\textbf{Odpowiedzi:}
\begin{description}
    \item [1:] a)$e^{i\pi}$, b)$\sqrt{2}e^{i\frac{\pi}{4}}$, c)$e^{i\frac{3}{2}\pi}=e^{-i\frac{1}{2}\pi}$, d)$2e^{-i\frac{1}{3}\pi}$, e)$\sqrt{53}e^{i(\arctan{(-\frac{7}{2})})}$
    \item [2:] a)$6e^{i\frac{3}{4}\pi}$, b)$\frac{2}{3}e^{-i\frac{1}{4}\pi}$, c)$2e^{i\frac{3}{4}\pi}$, d)$32e^{i\frac{5}{4}\pi}$, e)$2e^{-i\frac{1}{4}\pi}$, f)$\frac{1}{2}e^{-i\frac{1}{4}\pi}$
    \item [3:] a)$\{ 1+i, -1-i \}$, b)$\{ \frac{\sqrt{3}}{2} + \frac{1}{2}i, -\frac{\sqrt{3}}{2} + \frac{1}{2}i, -i \}$, c)$\{ 2i, -2i \}$, d)$\{ 1, i, -i, -1 \}$, e)$\{ \sqrt{3} + i, -\sqrt{3} + i, -2i \}$, f)$\{ \frac{\sqrt{3}+i}{2}, \frac{-\sqrt{3}-i}{2}, \frac{1 - i\sqrt{3}}{2} \}$, g)$\{ 1, \frac{1+i}{\sqrt{2}}, i, \frac{-1+i}{\sqrt{2}}, -1, \frac{-1-i}{\sqrt{2}}, -i, \frac{1-i}{\sqrt{2}} \}$
    \item [4:] a)$\{ -\frac{3}{2} + \frac{\sqrt{5}}{2}, -\frac{3}{2} - \frac{\sqrt{5}}{2} \}$, b)$\{ -2+i, -2-i \}$, c)$\{ \frac{\sqrt{2} - \sqrt{5}}{2} (1 + i), \frac{\sqrt{2} + \sqrt{5}}{2} (1 + i) \}$, d)$\{ 1+i\sqrt{3}, -2, 1-i\sqrt{3} \}$
\end{description}

\end{document}